%%%%%%%%%%%%%%%%%%%%%%%%%%%%%%%%%%%%%%%%%%
% Engineering problems / LaTeX Template
%		Semester 5
%		Institut d'Optique Graduate School
%%%%%%%%%%%%%%%%%%%%%%%%%%%%%%%%%%%%%%%%%%
%	5N-ONIP-Block1	/ Python for Science
%%%%%%%%%%%%%%%%%%%%%%%%%%%%%%%%%%%%%%%%%%
%
% Created by:
%	Julien VILLEMEJANE - 25/sep/2024
% Modified by:
%	
%
%%%%%%%%%%%%%%%%%%%%%%%%%%%%%%%%%%%%%%%%%%
% Professional Newsletter Template
% LaTeX Template
% Version 1.0 (09/03/14)
%
% Created by:
% Bob Kerstetter (https://www.tug.org/texshowcase/) and extensively modified by:
% Vel (vel@latextemplates.com)
% 
% This template has been downloaded from:
% http://www.LaTeXTemplates.com
%
% License:
% CC BY-NC-SA 3.0 (http://creativecommons.org/licenses/by-nc-sa/3.0/)
%
%%%%%%%%%%%%%%%%%%%%%%%%%%%%%%%%%%%%%%%%%

\documentclass[10pt]{article} % The default font size is 10pt; 11pt and 12pt are alternatives

%%%%%%%%%%%%%%%%%%%%%%%%%%%%%%%%%%%%%%%%%
% Professional Newsletter Template
% Structural Definitions File
% Version 1.0 (09/03/14)
%
% Created by:
% Vel (vel@latextemplates.com)
% 
% This file has been downloaded from:
% http://www.LaTeXTemplates.com
%
% License:
% CC BY-NC-SA 3.0 (http://creativecommons.org/licenses/by-nc-sa/3.0/)
%
%%%%%%%%%%%%%%%%%%%%%%%%%%%%%%%%%%%%%%%%%

%----------------------------------------------------------------------------------------
%	REQUIRED PACKAGES
%----------------------------------------------------------------------------------------

\usepackage{graphicx} % Required for including images
\usepackage{microtype} % Improved typography
\usepackage{multicol} % Used for the two-column layout of the document
\usepackage{booktabs} % Required for nice horizontal rules in tables
\usepackage{wrapfig} % Required for in-line images
\usepackage{float} % Required for forcing figures not to float with the [H] parameter

%------------------------------------------------
% Fonts

\usepackage{charter} % Use the Charter font as the main document font
\usepackage{courier} % Use the Courier font for \texttt (monospaced) only
\usepackage[T1]{fontenc} % Use T1 font encoding
\usepackage{lmodern}

%------------------------------------------------
% List Separation

\usepackage{enumitem} % Required to customize the list environments
\setlist{noitemsep,nolistsep} % Remove spacing before, after and within lists for a compact look

%------------------------------------------------
% Figure and Table Caption Styles

\usepackage{caption} % Required for changing caption styles
\captionsetup[table]{labelfont={bf,sf},labelsep=period,justification=justified} % Specify the table caption style
\captionsetup[figure]{labelfont={sf,bf},labelsep=period,justification=justified, font=small} % Specify the figure caption style
\setlength{\abovecaptionskip}{10pt} % Whitespace above captions

%------------------------------------------------
% Spacing Between Paragraphs

\makeatletter
\usepackage{parskip}
\setlength{\parskip}{6pt}
\newcommand{\@minipagerestore}{\setlength{\parskip}{6pt}}
\makeatother

%----------------------------------------------------------------------------------------
%	PAGE MARGINS AND SPACINGS
%----------------------------------------------------------------------------------------

\textwidth = 7 in % Text width
\textheight = 10 in % Text height
\oddsidemargin = -18pt % Left side margin on odd pages
\evensidemargin = -18pt % Left side margin on even pages
\topmargin = -36pt % Top margin
\headheight = 0pt % Remove the header by setting its space to 0
\headsep = 0pt % Remove the space between the header and top of the page
\parskip = 4pt % Space between paragraph
\parindent = 0.0in % Paragraph indentation
\pagestyle{empty} % Disable page numbering

%----------------------------------------------------------------------------------------
%	COLORS
%----------------------------------------------------------------------------------------

\usepackage[dvipsnames,svgnames]{xcolor} % Required to specify custom colors

\definecolor{altncolor}{rgb}{.8,0,0} % Dark red
%\definecolor{altncolor}{rgb}{.2,.4,.8} % Dark blue
%\definecolor{altncolor}{rgb}{.84,.16,.16} % Red

\usepackage[colorlinks=true, linkcolor=altncolor, anchorcolor=altncolor, citecolor=altncolor, filecolor=altncolor, menucolor=altncolor, urlcolor=altncolor]{hyperref} % Use the color defined above for all links

%----------------------------------------------------------------------------------------
%	DRAWINGS & CIRCUITS
%----------------------------------------------------------------------------------------

\usepackage{tikz}
\usepackage{pgfplots}

\usepackage[european, straightvoltages]{circuitikz}
\usetikzlibrary{babel, patterns, patterns.meta, calc}

%----------------------------------------------------------------------------------------
%	PDF
%----------------------------------------------------------------------------------------

\usepackage{pdfpages}

%----------------------------------------------------------------------------------------
%	BOX STYLES
%----------------------------------------------------------------------------------------

\usepackage[framemethod=TikZ]{mdframed}% Required for creating boxes
\mdfdefinestyle{sidebar}{
    linecolor=black, % Outer line color
    outerlinewidth=0.5pt, % Outer line width
    roundcorner=0pt, % Amount of corner rounding
    innertopmargin=10pt, % Top margin
    innerbottommargin=10pt, % Bottom margin
    innerrightmargin=10pt, % Right margin
    innerleftmargin=10pt, % Left margin
    backgroundcolor=white, % Box background color
    frametitlealignment=\centering,
    frametitlebackgroundcolor=gray!20, % Title background color
    frametitlerule=false, % Title rule - true or false
    frametitlerulecolor=white, % Title rule color
    frametitlerulewidth=0.5pt, % Title rule width
    frametitlefont=\Large\bfseries, % Title heading font specification
    font=\small
}

\mdfdefinestyle{intextbox}{
    linecolor=black, % Outer line color
    outerlinewidth=0.5pt, % Outer line width
    roundcorner=10pt, % Amount of corner rounding
    innertopmargin=7pt, % Top margin
    innerbottommargin=7pt, % Bottom margin
    innerrightmargin=7pt, % Right margin
    innerleftmargin=7pt, % Left margin
    backgroundcolor=white, % Box background color
    frametitlealignment=\centering,
    frametitlebackgroundcolor=gray!20, % Title background color
    frametitlerule=false, % Title rule - true or false
    frametitlerulecolor=white, % Title rule color
    frametitlerulewidth=0.5pt, % Title rule width
    frametitlefont=\Large\bfseries % Title heading font specification
}

\mdfdefinestyle{aavbox}{
    linecolor=black, % Outer line color
    outerlinewidth=0.2pt, % Outer line width
    roundcorner=5pt, % Amount of corner rounding
    innertopmargin=7pt, % Top margin
    innerbottommargin=7pt, % Bottom margin
    innerrightmargin=7pt, % Right margin
    innerleftmargin=7pt, % Left margin
    backgroundcolor=gray!10, % Box background color
    frametitlealignment=\centering,
    frametitlebackgroundcolor=gray!30, % Title background color
    frametitlerule=false, % Title rule - true or false
    frametitlerulecolor=white, % Title rule color
    frametitlerulewidth=0.2pt, % Title rule width
    frametitlefont=\Large\bfseries % Title heading font specification
}

%----------------------------------------------------------------------------------------
%	HEADING STYLE
%----------------------------------------------------------------------------------------

\newcommand{\heading}[2]{ % Define the \heading command
\vspace{#2} % White space above the heading
{\begin{center}\Large\textbf{#1}\end{center}} % The heading style
\vspace{#2} % White space below the heading
}

\newcommand{\BackToContents}{\hyperlink{contents}{{\small Back to Contents}}} % Define a command for linking back to the contents of the newsletter

\usepackage{listings}
\lstset{language = Python, 
	basicstyle={\small \color{black}}, 
	tabsize = 3,
	commentstyle=\color{black!70!green},
	linewidth=160mm,
	framexleftmargin=5mm, frame=shadowbox, rulesepcolor=\color{black},
  	numbers=left,
	xleftmargin=40pt,
	escapechar=|
} % Include the document which specifies all packages and structural customizations for this template
\usepackage{amsmath}

%----------------------------------------------------------------------------------------
%	DOCUMENT INFORMATIONS
%----------------------------------------------------------------------------------------
\def\module{Opto-Electronique}
\def\submodule{OptoElec}
\def\moduleSmall{5N-027-SCI / OE}
\def\year{2024-2025}
\def\problem{TD Systèmes et Signaux}
\def\problemName{OptoElec \& ONIP-1 / TD Systèmes et Signaux}

\def\validation{}

\def\scheduleCM{0}
\def\scheduleTD{0}
\def\scheduleTDcomputer{2}
\def\scheduleTP{0}

\def\workingTeam{Par binôme}

\def\workingSpecial{}

\def\keywords{Systèmes asservis;ALI;FFT}


\begin{document}
%----------------------------------------------------------------------------------------
%	HEADER IMAGE
%----------------------------------------------------------------------------------------

\begin{figure}[H]
\centering\includegraphics[width=0.3\linewidth]{../../../../../_assets/latex/logo_iogs.png}
\end{figure}

%----------------------------------------------------------------------------------------
%	MAIN BODY - FIRST PAGE
%----------------------------------------------------------------------------------------
%
\begin{minipage}[t]{1.0\linewidth} % Mini page taking up 65% of the actual page

\hypertarget{context}{\heading{\huge \problemName}{6pt}} % \hypertarget provides a label to reference using \hyperlink{label}{link text}

%-----------------------------------
\centerline {\rule{.70\linewidth}{.25pt}} % Horizontal line
\begin{center}
\textsc{\large Séance 2 / Signaux, images et FFT}
\end{center}
\centerline {\rule{.70\linewidth}{.25pt}} % Horizontal line
%-----------------------------------

\medskip

\begin{mdframed}[style=sidebar,frametitle={}]
\large
Un tutoriel sur la FFT est proposé sur la page suivante, section \textbf{\large FFT with Python} :

\begin{center}
\Large
\href{https://iogs-lense-training.github.io/python-for-science/}{https://iogs-lense-training.github.io/python-for-science/}
\end{center}

\end{mdframed}

%-----------------------------------
\textbf{Exercice 1 / FFT sur des signaux 1D}

\begin{enumerate}
	\item Définir une fonction qui génère un signal sinusoïdal d'une amplitude donnée et d'une fréquence donnée. Le vecteur temps devra également être passé en argument.
	\item Générer 2 signaux sinusoïdaux :
	\begin{itemize}
		\item le premier $sine\_a$ de fréquence $200\operatorname{Hz}$ et d'amplitude $1$
		\item le second $sine\_b$ de fréquence $287\operatorname{Hz}$ et d'amplitude $2$
	\end{itemize}
	\item Tracer ces deux signaux sur un même graphique.

	On se propose d'étudier un signal $sine\_c$ correspondant à la somme de ces deux signaux : $sine\_c = sine\_a + sine\_b$
	
	\item Tracer le signal $sine\_c$ sur un graphique.
	
	\medskip
	
	\item Calculer la FFT de ce signal et tracer cette réponse en fréquence sur un nouveau graphique, sans spécifier l'axe des fréquences. Que pouvez-vous conclure ?
	
	\medskip
	
	\item Construire l'axe des fréquences à l'aide de la fonction \textsl{numpy.fft.fftfreq}.
	\item Afficher la FFT du signal précédent en se basant cette fois-ci sur l'axe des fréquences générés à l'aide de \textsl{numpy.fft.fftfreq}.
\end{enumerate}



\centerline {\rule{.70\linewidth}{.25pt}} % Horizontal line
%-----------------------------------
\textbf{Exercice 2 / FFT sur des images}

On se propose d'étudier le script \textsl{image\_filter.py} (associé à l'image \textsc{test\_image.png}), disponible sur le site du LEnsE, rubrique Année / Première Année / Outils Numériques / TD Systèmes et Signaux.

\begin{enumerate}
	\item Tester ce script.
	\item Modifier la valeur du rayon (\textit{radius}) associée à la fonction \textsl{circular\_mask.py} et relancer le script.
	\item Expliquer les différentes étapes et l'impact sur l'image finale.
\end{enumerate}


\centerline {\rule{.70\linewidth}{.25pt}} % Horizontal line
%-----------------------------------
\textbf{Exercice 3 / FFT en 2D}

Le script \textsl{trame\_generator.py}, disponible sur le site du LEnsE, rubrique Année / Première Année / Outils Numériques / TD Systèmes et Signaux, contient une fonction permettant de générer des trames sinusoïdales en 2D.

\begin{enumerate}
	\item A l'aide de cette fonction, générer une trame sinusoïdale de taille $400\operatorname{pixels}$ par $300\operatorname{pixels}$ avec un angle de $63^\circ$ et un pas de $30\operatorname{pixels}$.
	\item Afficher l'image.
	\item Tracer et afficher la FFT de cette image (attention au \textit{shift}).
	\item Tester pour des pas et des angles différents.
\end{enumerate}

\end{minipage}

\end{document} 